\nonstopmode{}
\documentclass[a4paper]{book}
\usepackage[times,inconsolata,hyper]{Rd}
\usepackage{makeidx}
\makeatletter\@ifl@t@r\fmtversion{2018/04/01}{}{\usepackage[utf8]{inputenc}}\makeatother
% \usepackage{graphicx} % @USE GRAPHICX@
\makeindex{}
\begin{document}
\chapter*{}
\begin{center}
{\textbf{\huge Package `stacr'}}
\par\bigskip{\large \today}
\end{center}
\ifthenelse{\boolean{Rd@use@hyper}}{\hypersetup{pdftitle = {stacr: Tidy 'STAC' Workflows for R}}}{}
\ifthenelse{\boolean{Rd@use@hyper}}{\hypersetup{pdfauthor = {Chris Lyons}}}{}
\begin{description}
\raggedright{}
\item[Title]\AsIs{Tidy 'STAC' Workflows for R}
\item[Version]\AsIs{0.1.0}
\item[Description]\AsIs{Wraps the 'rstac' package with a pipe-friendly, tidy API. All results return 'tibbles' instead of nested lists. Ships with a catalog registry of known 'STAC' endpoints including Planetary Computer, Earth Search, and 'USGS', while supporting any 'STAC' API URL.}
\item[License]\AsIs{MIT + file LICENSE}
\item[URL]\AsIs{}\url{https://github.com/null-island-labs/stacr}\AsIs{}
\item[BugReports]\AsIs{}\url{https://github.com/null-island-labs/stacr/issues}\AsIs{}
\item[Depends]\AsIs{R (>= 4.1.0)}
\item[Imports]\AsIs{rstac, tibble, cli}
\item[Suggests]\AsIs{gdalcubes, jsonlite, leaflet, sf, testthat (>= 3.0.0), withr,
knitr, rmarkdown}
\item[Encoding]\AsIs{UTF-8}
\item[Roxygen]\AsIs{list(markdown = TRUE)}
\item[RoxygenNote]\AsIs{7.3.3}
\item[Config/testthat/edition]\AsIs{3}
\item[VignetteBuilder]\AsIs{knitr}
\item[NeedsCompilation]\AsIs{no}
\item[Author]\AsIs{Chris Lyons [aut, cre, cph]}
\item[Maintainer]\AsIs{Chris Lyons }\email{chrislyons@nullislandlabs.dev}\AsIs{}
\end{description}
\Rdcontents{Contents}
\HeaderA{stacr-package}{stacr: Tidy 'STAC' Workflows for R}{stacr.Rdash.package}
\aliasA{stacr}{stacr-package}{stacr}
\keyword{internal}{stacr-package}
%
\begin{Description}
Wraps the 'rstac' package with a pipe-friendly, tidy API. All results return 'tibbles' instead of nested lists. Ships with a catalog registry of known 'STAC' endpoints including Planetary Computer, Earth Search, and 'USGS', while supporting any 'STAC' API URL.
\end{Description}
%
\begin{Author}
\strong{Maintainer}: Chris Lyons \email{chrislyons@nullislandlabs.dev} [copyright holder]

\end{Author}
%
\begin{SeeAlso}
Useful links:
\begin{itemize}

\item{} \url{https://github.com/null-island-labs/stacr}
\item{} Report bugs at \url{https://github.com/null-island-labs/stacr/issues}

\end{itemize}


\end{SeeAlso}
\HeaderA{check\_suggested\_pkg}{Check that a suggested package is available}{check.Rul.suggested.Rul.pkg}
\keyword{internal}{check\_suggested\_pkg}
%
\begin{Description}
Check that a suggested package is available
\end{Description}
%
\begin{Usage}
\begin{verbatim}
check_suggested_pkg(pkg)
\end{verbatim}
\end{Usage}
%
\begin{Arguments}
\begin{ldescription}
\item[\code{pkg}] Package name as a string.
\end{ldescription}
\end{Arguments}
%
\begin{Value}
\code{TRUE} invisibly if available; throws an error otherwise.
\end{Value}
\HeaderA{has\_internet}{Check internet connectivity}{has.Rul.internet}
%
\begin{Description}
Tests whether a network connection to a STAC endpoint is available.
Used in \AsIs{\texttt{@examplesIf}} guards for network-dependent examples.
\end{Description}
%
\begin{Usage}
\begin{verbatim}
has_internet()
\end{verbatim}
\end{Usage}
%
\begin{Value}
\code{TRUE} if an internet connection is available, \code{FALSE} otherwise.
\end{Value}
%
\begin{Examples}
\begin{ExampleCode}
has_internet()
\end{ExampleCode}
\end{Examples}
\HeaderA{stac\_catalogs}{List Known STAC Catalogs}{stac.Rul.catalogs}
%
\begin{Description}
Returns a tibble of known STAC API endpoints bundled with the package.
Includes Planetary Computer, Earth Search, and USGS catalogs.
\end{Description}
%
\begin{Usage}
\begin{verbatim}
stac_catalogs()
\end{verbatim}
\end{Usage}
%
\begin{Value}
A \LinkA{tibble::tibble}{tibble::tibble} with columns:
\begin{description}

\item[name] Human-readable catalog name.
\item[url] Root URL of the STAC API endpoint.
\item[provider] Organization providing the catalog.

\end{description}

\end{Value}
%
\begin{Examples}
\begin{ExampleCode}
stac_catalogs()
\end{ExampleCode}
\end{Examples}
\HeaderA{stac\_collections}{List STAC Collections}{stac.Rul.collections}
%
\begin{Description}
Queries a STAC API endpoint and returns available collections as a tidy
tibble. Wraps \code{\LinkA{rstac::collections()}{rstac::collections()}} with tidy output.
\end{Description}
%
\begin{Usage}
\begin{verbatim}
stac_collections(url)
\end{verbatim}
\end{Usage}
%
\begin{Arguments}
\begin{ldescription}
\item[\code{url}] Character. Root URL of a STAC API endpoint
(e.g., \code{"https://earth-search.aws.element84.com/v1"}).
\end{ldescription}
\end{Arguments}
%
\begin{Value}
A \LinkA{tibble::tibble}{tibble::tibble} with one row per collection and columns:
\begin{description}

\item[id] Collection identifier.
\item[title] Human-readable title.
\item[description] Collection description.

\end{description}

\end{Value}
%
\begin{Examples}
\begin{ExampleCode}

stac_collections("https://earth-search.aws.element84.com/v1")

\end{ExampleCode}
\end{Examples}
\HeaderA{stac\_download}{Download STAC Assets}{stac.Rul.download}
%
\begin{Description}
Downloads assets from STAC items returned by \code{\LinkA{stac\_search()}{stac.Rul.search}} or
\code{\LinkA{stac\_items()}{stac.Rul.items}}. Wraps \code{\LinkA{rstac::assets\_download()}{rstac::assets.Rul.download()}} with tidy output.
\end{Description}
%
\begin{Usage}
\begin{verbatim}
stac_download(items, assets = NULL, output_dir = tempdir(), overwrite = FALSE)
\end{verbatim}
\end{Usage}
%
\begin{Arguments}
\begin{ldescription}
\item[\code{items}] An rstac \code{doc\_items} object as returned by \code{\LinkA{stac\_search()}{stac.Rul.search}} or
\code{\LinkA{stac\_items()}{stac.Rul.items}} with the \code{raw} attribute, or the raw rstac result directly.

\item[\code{assets}] Character vector. Asset names to download (e.g.,
\code{c("red", "green", "blue")}). If \code{NULL}, downloads all assets.

\item[\code{output\_dir}] Character. Directory where files are saved.
Defaults to \code{\LinkA{tempdir()}{tempdir}}.

\item[\code{overwrite}] Logical. Overwrite existing files? Defaults to \code{FALSE}.
\end{ldescription}
\end{Arguments}
%
\begin{Value}
The rstac items object (invisibly) with assets downloaded to
\code{output\_dir}. Downloaded file paths can be retrieved with
\code{\LinkA{rstac::assets\_url()}{rstac::assets.Rul.url()}}.
\end{Value}
%
\begin{Examples}
\begin{ExampleCode}


# Search for items first
items <- stac_search(
  url = "https://earth-search.aws.element84.com/v1",
  collections = "sentinel-2-l2a",
  bbox = c(-84.5, 38.0, -84.3, 38.2),
  limit = 1
)


\end{ExampleCode}
\end{Examples}
\HeaderA{stac\_items}{List Items in a STAC Collection}{stac.Rul.items}
%
\begin{Description}
Retrieves items from a specific collection in a STAC API endpoint and
returns them as a tidy tibble. Wraps \code{\LinkA{rstac::items()}{rstac::items()}} with tidy output.
\end{Description}
%
\begin{Usage}
\begin{verbatim}
stac_items(url, collection, limit = 100L)
\end{verbatim}
\end{Usage}
%
\begin{Arguments}
\begin{ldescription}
\item[\code{url}] Character. Root URL of a STAC API endpoint
(e.g., \code{"https://earth-search.aws.element84.com/v1"}).

\item[\code{collection}] Character. The collection ID to list items from.

\item[\code{limit}] Integer. Maximum number of items to return. Defaults to 100.
\end{ldescription}
\end{Arguments}
%
\begin{Value}
A \LinkA{tibble::tibble}{tibble::tibble} with one row per item and columns:
\begin{description}

\item[id] Item identifier.
\item[collection] Collection the item belongs to.
\item[datetime] Acquisition datetime as a character string.
\item[bbox] Bounding box as a numeric list column.
\item[geometry] GeoJSON geometry as a list column.
\item[assets] Character vector of available asset names.

\end{description}

\end{Value}
%
\begin{Examples}
\begin{ExampleCode}


stac_items(
  url = "https://earth-search.aws.element84.com/v1",
  collection = "sentinel-2-l2a",
  limit = 5
)


\end{ExampleCode}
\end{Examples}
\HeaderA{stac\_map}{Map STAC Item Footprints}{stac.Rul.map}
%
\begin{Description}
Creates an interactive 'leaflet' map showing the spatial footprints of
STAC items returned by \code{\LinkA{stac\_search()}{stac.Rul.search}} or \code{\LinkA{stac\_items()}{stac.Rul.items}}. Requires the
'leaflet' and 'sf' packages to be installed.
\end{Description}
%
\begin{Usage}
\begin{verbatim}
stac_map(items)
\end{verbatim}
\end{Usage}
%
\begin{Arguments}
\begin{ldescription}
\item[\code{items}] A \LinkA{tibble::tibble}{tibble::tibble} of STAC items as returned by
\code{\LinkA{stac\_search()}{stac.Rul.search}} or \code{\LinkA{stac\_items()}{stac.Rul.items}}, with a \code{geometry} list column.
\end{ldescription}
\end{Arguments}
%
\begin{Value}
A \code{leaflet} htmlwidget object.
\end{Value}
%
\begin{Examples}
\begin{ExampleCode}


items <- stac_search(
  url = "https://earth-search.aws.element84.com/v1",
  collections = "sentinel-2-l2a",
  bbox = c(-84.5, 38.0, -84.3, 38.2),
  limit = 5
)
stac_map(items)


\end{ExampleCode}
\end{Examples}
\HeaderA{stac\_registry}{Built-in catalog registry}{stac.Rul.registry}
\keyword{internal}{stac\_registry}
%
\begin{Description}
Returns a tibble of known STAC API endpoints bundled with the package.
\end{Description}
%
\begin{Usage}
\begin{verbatim}
stac_registry()
\end{verbatim}
\end{Usage}
%
\begin{Value}
A \LinkA{tibble::tibble}{tibble::tibble} with columns: \code{name}, \code{url}, \code{provider}.
\end{Value}
\HeaderA{stac\_search}{Search a STAC API}{stac.Rul.search}
%
\begin{Description}
Searches a STAC API endpoint for items matching the given filters and
returns results as a tidy tibble. Wraps \code{\LinkA{rstac::stac\_search()}{rstac::stac.Rul.search()}} with
tidy output.
\end{Description}
%
\begin{Usage}
\begin{verbatim}
stac_search(
  url,
  collections = NULL,
  bbox = NULL,
  datetime = NULL,
  limit = 100L
)
\end{verbatim}
\end{Usage}
%
\begin{Arguments}
\begin{ldescription}
\item[\code{url}] Character. Root URL of a STAC API endpoint
(e.g., \code{"https://earth-search.aws.element84.com/v1"}).

\item[\code{collections}] Character vector. Collection IDs to search within.

\item[\code{bbox}] Numeric vector of length 4: \code{c(xmin, ymin, xmax, ymax)}.
Only items intersecting this bounding box are returned.

\item[\code{datetime}] Character. Date/time filter as a single datetime or
range (e.g., \code{"2024-01-01/2024-12-31"}). Follows RFC 3339.

\item[\code{limit}] Integer. Maximum number of items to return. Defaults to 100.
\end{ldescription}
\end{Arguments}
%
\begin{Value}
A \LinkA{tibble::tibble}{tibble::tibble} with one row per item and columns:
\begin{description}

\item[id] Item identifier.
\item[collection] Collection the item belongs to.
\item[datetime] Acquisition datetime as a character string.
\item[bbox] Bounding box as a numeric list column.
\item[geometry] GeoJSON geometry as a list column.
\item[assets] Character vector of available asset names.

\end{description}

\end{Value}
%
\begin{Examples}
\begin{ExampleCode}


stac_search(
  url = "https://earth-search.aws.element84.com/v1",
  collections = "sentinel-2-l2a",
  bbox = c(-84.5, 38.0, -84.3, 38.2),
  limit = 5
)


\end{ExampleCode}
\end{Examples}
\HeaderA{stac\_search\_raw}{Search STAC and Return Raw rstac Result}{stac.Rul.search.Rul.raw}
%
\begin{Description}
Like \code{\LinkA{stac\_search()}{stac.Rul.search}} but returns the raw rstac \code{doc\_items} object
instead of a tibble. Useful as input for \code{\LinkA{stac\_download()}{stac.Rul.download}} and
\code{\LinkA{stac\_to\_cube()}{stac.Rul.to.Rul.cube}}.
\end{Description}
%
\begin{Usage}
\begin{verbatim}
stac_search_raw(
  url,
  collections = NULL,
  bbox = NULL,
  datetime = NULL,
  limit = 100L
)
\end{verbatim}
\end{Usage}
%
\begin{Arguments}
\begin{ldescription}
\item[\code{url}] Character. Root URL of a STAC API endpoint
(e.g., \code{"https://earth-search.aws.element84.com/v1"}).

\item[\code{collections}] Character vector. Collection IDs to search within.

\item[\code{bbox}] Numeric vector of length 4: \code{c(xmin, ymin, xmax, ymax)}.
Only items intersecting this bounding box are returned.

\item[\code{datetime}] Character. Date/time filter as a single datetime or
range (e.g., \code{"2024-01-01/2024-12-31"}). Follows RFC 3339.

\item[\code{limit}] Integer. Maximum number of items to return. Defaults to 100.
\end{ldescription}
\end{Arguments}
%
\begin{Value}
An rstac \code{doc\_items} object.
\end{Value}
%
\begin{Examples}
\begin{ExampleCode}


raw <- stac_search_raw(
  url = "https://earth-search.aws.element84.com/v1",
  collections = "sentinel-2-l2a",
  bbox = c(-84.5, 38.0, -84.3, 38.2),
  limit = 1
)


\end{ExampleCode}
\end{Examples}
\HeaderA{stac\_to\_cube}{Convert STAC Items to a gdalcubes Image Collection}{stac.Rul.to.Rul.cube}
%
\begin{Description}
Bridges STAC search results to 'gdalcubes' for raster data cube analysis.
Requires the 'gdalcubes' package to be installed.
\end{Description}
%
\begin{Usage}
\begin{verbatim}
stac_to_cube(items, asset_names = NULL, ...)
\end{verbatim}
\end{Usage}
%
\begin{Arguments}
\begin{ldescription}
\item[\code{items}] An rstac \code{doc\_items} object as returned by \code{\LinkA{stac\_search\_raw()}{stac.Rul.search.Rul.raw}}
or \code{\LinkA{rstac::post\_request()}{rstac::post.Rul.request()}}.

\item[\code{asset\_names}] Character vector. Asset names to include in the
image collection (e.g., \code{c("red", "green", "blue")}).

\item[\code{...}] Additional arguments passed to
\code{gdalcubes::stac\_image\_collection()}.
\end{ldescription}
\end{Arguments}
%
\begin{Value}
A \code{gdalcubes} image collection object.
\end{Value}
%
\begin{Examples}
\begin{ExampleCode}


raw <- stac_search_raw(
  url = "https://earth-search.aws.element84.com/v1",
  collections = "sentinel-2-l2a",
  bbox = c(-84.5, 38.0, -84.3, 38.2),
  limit = 5
)
cube <- stac_to_cube(raw, asset_names = c("red", "green", "blue"))


\end{ExampleCode}
\end{Examples}
\HeaderA{tidy\_collections}{Convert STAC collections to a tidy tibble}{tidy.Rul.collections}
\keyword{internal}{tidy\_collections}
%
\begin{Description}
Convert STAC collections to a tidy tibble
\end{Description}
%
\begin{Usage}
\begin{verbatim}
tidy_collections(collections)
\end{verbatim}
\end{Usage}
%
\begin{Arguments}
\begin{ldescription}
\item[\code{collections}] An rstac \code{doc\_collections} object.
\end{ldescription}
\end{Arguments}
%
\begin{Value}
A \LinkA{tibble::tibble}{tibble::tibble} with columns: \code{id}, \code{title}, \code{description}.
\end{Value}
\HeaderA{tidy\_items}{Convert STAC items to a tidy tibble}{tidy.Rul.items}
\keyword{internal}{tidy\_items}
%
\begin{Description}
Extracts key fields from an rstac items response (\code{doc\_items}) and returns
a flat tibble with one row per item.
\end{Description}
%
\begin{Usage}
\begin{verbatim}
tidy_items(items)
\end{verbatim}
\end{Usage}
%
\begin{Arguments}
\begin{ldescription}
\item[\code{items}] An rstac \code{doc\_items} object (result of \code{\LinkA{rstac::post\_request()}{rstac::post.Rul.request()}}
or \code{\LinkA{rstac::get\_request()}{rstac::get.Rul.request()}} on a search or items endpoint).
\end{ldescription}
\end{Arguments}
%
\begin{Value}
A \LinkA{tibble::tibble}{tibble::tibble} with columns: \code{id}, \code{collection}, \code{datetime},
\code{bbox}, \code{geometry}, and \code{assets}. Additional property columns are included
from \code{rstac::items\_as\_tibble()}.
\end{Value}
\HeaderA{validate\_bbox}{Validate a bounding box}{validate.Rul.bbox}
\keyword{internal}{validate\_bbox}
%
\begin{Description}
Checks that \code{bbox} is a numeric vector of length 4 with valid coordinate
ordering: \code{c(xmin, ymin, xmax, ymax)}.
\end{Description}
%
\begin{Usage}
\begin{verbatim}
validate_bbox(bbox)
\end{verbatim}
\end{Usage}
%
\begin{Arguments}
\begin{ldescription}
\item[\code{bbox}] Numeric vector of length 4.
\end{ldescription}
\end{Arguments}
%
\begin{Value}
\code{bbox} invisibly if valid; throws an error otherwise.
\end{Value}
\HeaderA{validate\_url}{Validate a STAC endpoint URL}{validate.Rul.url}
\keyword{internal}{validate\_url}
%
\begin{Description}
Validate a STAC endpoint URL
\end{Description}
%
\begin{Usage}
\begin{verbatim}
validate_url(url)
\end{verbatim}
\end{Usage}
%
\begin{Arguments}
\begin{ldescription}
\item[\code{url}] Character string.
\end{ldescription}
\end{Arguments}
%
\begin{Value}
\code{url} invisibly if valid; throws an error otherwise.
\end{Value}
\printindex{}
\end{document}
